\section{Evaluation Details}
\label{app:evaluation}

This appendix fixes the collection rules behind the four result figures. The
release gate accepts a figure only when its CSV, provenance record, inputs, and
rendered output agree by SHA-256 digest.

\subsection{Extension Tasks and Oracles}
\label{app:extension}

The 24 evaluation tasks and the demonstration bank contain disjoint feature
identities. Demonstrations are nested prefixes of one deterministic ranking,
so increasing the count adds context without replacing earlier examples. A
sealed reference bundle is unavailable during generation and opens only after
all model outputs are frozen. Each output is classified at its first failing
phase: parse, type, build, or semantic oracle. Formatting changes and supplied
harness boilerplate do not count toward generated target tokens. Every row
binds the model revision, task specification, API card, prompt, output,
demonstration IDs, and sampling seed.

\subsection{Patch Reduction and Ownership Zones}
\label{app:patch}

Each matched implementation starts at a pinned clean revision. Hunk-level
delta debugging visits candidate hunks in a fixed order, removes one, and keeps
the removal only when the build and semantic oracle still pass. Iteration ends
at a 1-minimal fixed point under that partition. The frozen counting policy
excludes generated files, vendored sources, lock files, and formatter-only
changes. Implementation and test changes remain separate. Ownership zones are
source packages or build targets with one public responsibility; crossed edges
count dependencies between those zones. The artifact binds the base revision,
final patch, reduction trace, zone map, and oracle output.

\subsection{Update Adapter Contract}
\label{app:update}

Each adapter consumes the same neutral typed graph and exposes five logical
stages: analysis, canonicalization, target selection, memory planning, and
artifact-manifest construction. Counters report visited graph entities and
executed stages over the timed interval. A case selects an early, middle, or
late operation, changes either operation metadata or a result type, and targets
either its affected cone or an unrelated entity. The full and update paths
start from identical published state. Stage digests and the final canonical
digest must match within a system and across systems. The validator requires
every cell in the 15-by-3-by-3-by-2-by-2-by-2-by-100 matrix.

\subsection{Artifact Performance Protocol}
\label{app:performance}

The operator corpus contains four cases in each of six families. Model and
operator variants consume byte-identical, hash-bound tensors. Collection pins
compiler and runtime revisions, release flags, CPU affinity, thread count, and
seed. Timing begins immediately before artifact invocation and ends after
completion; graph loading, input generation, and compilation stay outside the
interval. Each supported pair runs ten warm-ups and 100 recorded iterations.
Floating-point outputs use dtype-specific absolute and relative tolerances;
integer and Boolean outputs require exact equality. Unsupported pairs retain a
reason and contribute to coverage, but not to latency ratios.
